The last two lines of Algorithm~\ref{alg:adaptive} apply a reweighting step to the
parameter produced by the concentration steps. We study the step as a map on its own,
applied to a fixed parameter $\theta=(\theta_1,\dots,\theta_K)\in\R^{Kd}$, and let the
sample size grow. The step, as implemented, is the following. For unit $i$ let
$z_i(\theta)=\argmin_k|y_i-x_i^{\top}\theta_k|$ be its nearest surface, the smallest index in
case of a tie, and $r_i(\theta)=\min_k|y_i-x_i^{\top}\theta_k|$ its absolute residual from
that surface. Let $n_k=\#\{i:z_i(\theta)=k\}$. The scale of group $k$ is
$\hat s_k=1.4826\,\hat m_k$, where $\hat m_k$ is a sample median of
$\{r_i(\theta):z_i(\theta)=k\}$, by which we mean any value between the two middle order
statistics of that set; the implementation takes the square root of the median of the
squared residuals, which lies between them. When $n_k<10$ the median over all $n$ residuals
is used instead. The flagged set and the flagged fraction are
\[
\mathcal F_n(\theta)=\{i:\ r_i(\theta)>c\,\hat s_{z_i(\theta)}\},\qquad
\hat\alpha_n(\theta)=|\mathcal F_n(\theta)|/n .
\]
The refitted parameter $\hat\theta_k$ is the least-squares fit of $y$ on $x$ over
$\{i:z_i(\theta)=k,\ i\notin\mathcal F_n(\theta)\}$ when this set has at least $d$ units, and
$\hat\theta_k=\theta_k$ otherwise.

The population counterparts are defined for a probability measure $P$ on
$\mathcal U=\R^{d}\times\R$ and a unit $u=(x,y)$:
\begin{align*}
h_\theta(u)&=\argmin_k|y-x^{\top}\theta_k|,\qquad
r_\theta(u)=\min_k|y-x^{\top}\theta_k|,\qquad
A_k(\theta)=\{u:h_\theta(u)=k\},\\
m_k(\theta)&=\inf\{t\ge0:\ P(r_\theta\le t\mid A_k(\theta))\ge\tfrac12\},\qquad
s_k(\theta)=1.4826\,m_k(\theta),\\
W_k(\theta)&=A_k(\theta)\cap\{r_\theta\le c\,s_k(\theta)\},\qquad
\alpha(\theta)=P\bigl(r_\theta>c\,s_{h_\theta}(\theta)\bigr)=1-\textstyle\sum_kP(W_k(\theta)),\\
M_k(\theta)&=\E\bigl[xx^{\top}\one_{W_k(\theta)}\bigr],\qquad
v_k(\theta)=\E\bigl[xy\,\one_{W_k(\theta)}\bigr],\qquad
\theta^{\infty}_k(\theta)=M_k(\theta)^{-1}v_k(\theta).
\end{align*}
Here $m_k(\theta)$ is the median of the residual $r_\theta$ among the units assigned to
surface $k$, $s_k(\theta)$ the corresponding scale, $W_k(\theta)$ the set of units assigned to
$k$ and not flagged, $\alpha(\theta)$ the population flagged fraction, and
$\theta^{\infty}(\theta)$ the population refit. The dependence on $\theta$ is dropped when
$\theta$ is fixed.

\begin{assumption}[Clean law]
\label{ass:clean}
$P=(1-\varepsilon)P_0+\varepsilon C$ with $\varepsilon\in[0,1)$. Under $P_0$ a unit is
generated by a label $z\in\{1,\dots,K\}$ with $\Pr(z=k)=\pi_k>0$, a covariate $x$ whose
conditional law given $z=k$ has $\E_0[\|x\|^{2}]<\infty$, and an error $e\sim\mathcal N(0,1)$
independent of $(x,z)$, through $y=x^{\top}\theta^{0}_z+\sigma_ze$ with $\sigma_k>0$. The
contaminating law $C$ is arbitrary.
\end{assumption}

\begin{assumption}[Contamination beyond the cut-off]
\label{ass:far}
At the given $\theta$, there is $c'>c$ such that
$C\bigl(r_\theta\ge c'\,s_{h_\theta}(\theta)\bigr)=1$: every contaminated unit lies beyond
$c'$ scales from every surface of $\theta$, the scale being that of the surface nearest to it.
\end{assumption}

The scale $s_k(\theta)$ in Assumption~\ref{ass:far} is computed under $P$, so the
assumption is a condition on $P$; it is the form of ``beyond the cut-off'' that refers to the
scale the step actually uses. In the designs of the study the gross outliers are shifted by $15$
to $30$ noise scales from their line (moderate outliers by $4$ to $8$), and uniform noise is
redrawn until it lies more than three noise scales from both true lines, so at $\theta^{0}$ the
assumption holds for all or nearly all contaminated units whenever $s_k(\theta^{0})$ is close to
$\sigma_k$; a shifted unit can still fall near another group's line, and the assumption then
fails for that unit.

\begin{assumption}[Regularity at $\theta$]
\label{ass:reg}
For every $k$: $P(A_k(\theta))>0$; the median of $r_\theta$ on $A_k(\theta)$ is positive and
unique, that is, $m_k(\theta)>0$, $P(r_\theta\le t\mid A_k)<\tfrac12$ for $t<m_k$ and
$>\tfrac12$ for $t>m_k$; and $M_k(\theta)$ is nonsingular.
\end{assumption}

\begin{lemma}[Plug-in cut-off]
\label{lem:plugin}
Let $u_1,u_2,\dots$ be independent with law $P$, let $k$ be fixed, let $\hat s_k$ be random
with $\hat s_k\to s_k$ almost surely, and let $g:\mathcal U\to\R^{p}$ be measurable with
$P(A_k,\,r_\theta=cs_k)=0$ and $\E[\|g\|\one\{A_k,\,r_\theta\le cs_k+\delta_0\}]<\infty$ for
some $\delta_0>0$. Then, almost surely,
\[
\frac1n\sum_{i=1}^{n}g(u_i)\one\{z_i(\theta)=k,\ r_i(\theta)\le c\hat s_k\}
\;\to\;\E\bigl[g\,\one_{W_k}\bigr].
\]
\end{lemma}

\begin{proof}
Fix $\delta\in(0,\delta_0)$. Almost surely $|c\hat s_k-cs_k|\le\delta$ for all large $n$, and
then the average on the left differs from
$n^{-1}\sum_ig(u_i)\one\{z_i=k,\,r_i\le cs_k\}$ by at most
$n^{-1}\sum_i\|g(u_i)\|\one\{z_i=k,\,cs_k-\delta<r_i\le cs_k+\delta\}$ in norm. By the strong
law of large numbers the second average converges to $\E[g\one_{W_k}]$ and the bound to
$\E[\|g\|\one\{A_k,\,cs_k-\delta<r_\theta\le cs_k+\delta\}]$, both integrable by hypothesis.
Letting $\delta\downarrow0$ along a sequence, the last expectation tends to
$\E[\|g\|\one\{A_k,\,r_\theta=cs_k\}]=0$ by dominated convergence.
\end{proof}

\begin{theorem}[One reweighting step at a fixed input]
\label{thm:stepfixed}
Let $u_1,\dots,u_n$ be independent with law $P$, fix $\theta$, and let
Assumptions~\ref{ass:clean}--\ref{ass:reg} hold at $\theta$. Then, almost surely as
$n\to\infty$:
\begin{enumerate}[label=(\roman*),leftmargin=2.2em,itemsep=2pt]
\item $\hat s_k\to s_k(\theta)$ for every $k$;
\item $\hat\alpha_n(\theta)\to\alpha(\theta)=\varepsilon+(1-\varepsilon)\kappa(\theta)$,
where $\kappa(\theta)=P_0\bigl(r_\theta>c\,s_{h_\theta}(\theta)\bigr)$ is the flagged
fraction among clean units;
\item $\hat\theta\to\theta^{\infty}(\theta)$, and
$\theta^{\infty}_k(\theta)=\E_0[xx^{\top}\one_{W_k}]^{-1}\E_0[xy\,\one_{W_k}]$: the limit of
the refit is an average over the clean law alone, and the contamination enters it only through
the scales $s_k(\theta)$ that define $W_k$.
\end{enumerate}
\end{theorem}

\begin{proof}
\emph{Step 1: no mass at the cut-off.} Under $P_0$, $y$ has a density conditionally on
$(x,z)$, so $P_0(|y-x^{\top}\theta_j|=t)=0$ for every $j$ and $t$, and
$P_0(r_\theta=t)\le\sum_jP_0(|y-x^{\top}\theta_j|=t)=0$. Under $C$,
Assumption~\ref{ass:far} gives $C(A_k,\,r_\theta<c's_k)=0$. Hence
$P(A_k,\,r_\theta=cs_k)=0$, and for $\delta_0=\tfrac12(c'-c)\min_ks_k$, which is positive by
Assumption~\ref{ass:reg}, and every $g\ge0$,
\begin{equation}
\E\bigl[g\,\one\{A_k,\,r_\theta\le cs_k+\delta_0\}\bigr]
=(1-\varepsilon)\E_0\bigl[g\,\one\{A_k,\,r_\theta\le cs_k+\delta_0\}\bigr].
\label{eq:adp-cleanonly}
\end{equation}
In particular $P(W_k)=(1-\varepsilon)P_0(W_k)$, $M_k=(1-\varepsilon)\E_0[xx^{\top}\one_{W_k}]$
and $v_k=(1-\varepsilon)\E_0[xy\one_{W_k}]$; the factors $1-\varepsilon$ cancel in
$\theta^{\infty}_k$, which proves the second claim of (iii).

\emph{Step 2: scales.} By the strong law, $n_k/n\to P(A_k)>0$, so $n_k\ge10$ for all large
$n$ and the pooled fallback is inactive. For $t\ge0$ put
$\hat F_k(t)=n_k^{-1}\#\{i:z_i=k,\,r_i\le t\}$ and $F_k(t)=P(r_\theta\le t\mid A_k)$; then
$\hat F_k(t)\to F_k(t)$ almost surely for each fixed $t$, as a ratio of two averages. Fix
$\delta>0$. By the uniqueness of the median, $F_k(m_k-\delta)<\tfrac12<F_k(m_k+\delta)$, so
for all large $n$, $\hat F_k(m_k-\delta)<\tfrac12<\hat F_k(m_k+\delta)$: fewer than half of
the residuals of group $k$ are at most $m_k-\delta$ and more than half are at most
$m_k+\delta$. Both middle order statistics, and hence $\hat m_k$, then lie in
$(m_k-\delta,m_k+\delta]$. Taking $\delta=1/j$, $j\ge1$, gives $\hat m_k\to m_k$ almost
surely, which is (i).

\emph{Step 3: flagged fraction.} Apply Lemma~\ref{lem:plugin} with $g\equiv1$, the
hypotheses holding by Step 1 and (i): $n^{-1}\#\{i:z_i=k,\,r_i\le c\hat s_k\}\to P(W_k)$ for
each $k$. Summing over $k$, $1-\hat\alpha_n(\theta)\to\sum_kP(W_k)=1-\alpha(\theta)$. By
Step 1, $\sum_kP(W_k)=(1-\varepsilon)\sum_kP_0(W_k)=(1-\varepsilon)(1-\kappa(\theta))$, which
rearranges to (ii).

\emph{Step 4: refit.} Apply Lemma~\ref{lem:plugin} with $g=xx^{\top}$ and with $g=xy$;
the integrability hypothesis holds by \eqref{eq:adp-cleanonly}, since
$\E_0\|x\|^{2}<\infty$ and $\E_0y^{2}<\infty$ under Assumption~\ref{ass:clean}. Thus
$n^{-1}\sum_ix_ix_i^{\top}\one\{z_i=k,\,i\notin\mathcal F_n\}\to M_k$ and
$n^{-1}\sum_ix_iy_i\one\{z_i=k,\,i\notin\mathcal F_n\}\to v_k$ almost surely. Since $M_k$ is
nonsingular, $P(W_k)>0$, so the number of unflagged units of group $k$ exceeds $d$ for all
large $n$, the fallback is inactive, and $\hat\theta_k$ is the image of the two averages
under the map $(M,v)\mapsto M^{-1}v$, which is continuous at $(M_k,v_k)$. This is (iii).
\end{proof}

Theorem~\ref{thm:stepfixed} identifies the limit of the step at any input satisfying the
assumptions. It does not say that the input of the step in Algorithm~\ref{alg:adaptive} is
close to $\theta^{0}$. The next theorem evaluates the limits at $\theta=\theta^{0}$, where they
can be made explicit, and quantifies the effect of overlap between the groups. Write
\[
\eta_k=P_0\bigl(h_{\theta^{0}}\ne k\bigm|z=k\bigr),\qquad
\eta=\sum_k\pi_k\eta_k=P_0\bigl(h_{\theta^{0}}\ne z\bigr),
\]
for the probability that a clean unit of group $k$, or a clean unit, is nearer to another
true surface than to its own. Let $p_k=P(A_k(\theta^{0}))$ and
\[
\delta_k=1-\frac{(1-\varepsilon)\pi_k(1-\eta_k)}{p_k},\qquad \nu_k=\delta_k+\tfrac12\eta_k,
\]
so that $\delta_k$ is the fraction of the units assigned to surface $k$ that are not clean
units of group $k$; it satisfies $\delta_kp_k\le(1-\varepsilon)\eta+\varepsilon$, since the
clean units of other groups assigned to $k$ have mass at most $(1-\varepsilon)\eta$. Finally
let
\[
q(u)=1.4826\,\Phi^{-1}\bigl(\tfrac{1+u}{2}\bigr),\qquad 0<u<1,
\]
an increasing function with $q(\tfrac12)=1.4826\,\Phi^{-1}(\tfrac34)=1.0000$ to four
decimals; $q(u)\sigma$ is the scale the step assigns to a $\mathcal N(0,\sigma^{2})$ residual
whose $u$-quantile in absolute value is used in place of the median.

\begin{theorem}[The step at the true parameter]
\label{thm:steptrue}
Let Assumption~\ref{ass:clean} hold, let $\theta=\theta^{0}$, and suppose that
$P_0(x^{\top}(\theta^{0}_k-\theta^{0}_j)=0)=0$ for $j\ne k$ (no two surfaces coincide on a
set of covariates of positive mass), that $\nu_k<\tfrac12$ for every $k$, and that
Assumption~\ref{ass:far} holds at $\theta^{0}$.
\begin{enumerate}[label=(\roman*),leftmargin=2.2em,itemsep=2pt]
\item For every $k$, $p_k>0$ and the median of $r_{\theta^{0}}$ on $A_k(\theta^{0})$ is
positive and unique, and
\[
q(\tfrac12-\nu_k)\;\le\;\frac{s_k(\theta^{0})}{\sigma_k}\;\le\;q(\tfrac12+\nu_k).
\]
\item The clean flagged fraction satisfies
\[
\sum_k\pi_k\,2\Phi\bigl(-c\,q(\tfrac12+\nu_k)\bigr)-\eta
\;\le\;\kappa(\theta^{0})\;\le\;
\sum_k\pi_k\,2\Phi\bigl(-c\,q(\tfrac12-\nu_k)\bigr)+\eta,
\]
and $\alpha(\theta^{0})=\varepsilon+(1-\varepsilon)\kappa(\theta^{0})$.
\item Suppose in addition that $\|x\|\le M_x$ $P_0$-almost surely and that
\[
\lambda_k:=\bigl\{2\Phi\bigl(c\,q(\tfrac12-\nu_k)\bigr)-1\bigr\}\,
\lambda_{\min}\bigl(\E_0[xx^{\top}\mid z=k]\bigr)-M_x^{2}\eta_k\;>\;0 .
\]
Then $M_k(\theta^{0})$ is nonsingular with $\lambda_{\min}(M_k(\theta^{0}))\ge(1-\varepsilon)\pi_k\lambda_k$,
Assumption~\ref{ass:reg} holds at $\theta^{0}$, Theorem~\ref{thm:stepfixed} applies, and
\[
\bigl\|\theta^{\infty}_k(\theta^{0})-\theta^{0}_k\bigr\|\;\le\;
\frac{2\,c\,s_k(\theta^{0})\,M_x}{\pi_k\lambda_k}\;\eta .
\]
\end{enumerate}
In particular, if $\varepsilon=0$, $\eta=0$ and every $\E_0[xx^{\top}\mid z=k]$ is nonsingular, as in (iii), then $s_k(\theta^{0})=q(\tfrac12)\sigma_k$,
$\kappa(\theta^{0})=2\Phi(-c\,q(\tfrac12))$ and $\theta^{\infty}(\theta^{0})=\theta^{0}$:
the true parameter is a fixed point of the population reweighting step, and the flagged
fraction is $2\Phi(-c)$ up to the rounding in $1.4826$.
\end{theorem}

\begin{proof}
Throughout, $h=h_{\theta^{0}}$, $r=r_{\theta^{0}}$, $A_k=A_k(\theta^{0})$,
$s_k=s_k(\theta^{0})$ and $G_k(t)=\Pr(\sigma_k|e|\le t)=2\Phi(t/\sigma_k)-1$.

\emph{Geometry of the assignment.} For $j\ne k$ put $\Delta_j(x)=x^{\top}(\theta^{0}_k-\theta^{0}_j)$,
which is nonzero $P_0$-almost surely. On $\{z=k\}$ we have $y-x^{\top}\theta^{0}_k=\sigma_ke$
and $y-x^{\top}\theta^{0}_j=\sigma_ke+\Delta_j$, and $|a|<|a+\Delta|$ if and only if
$\Delta(2a+\Delta)>0$. Hence, given $x$, the event $\{h=k\}$ is $\{-L(x)<\sigma_ke<U(x)\}$
with $L(x)=\min\{\Delta_j/2:\Delta_j>0\}$ and $U(x)=\min\{-\Delta_j/2:\Delta_j<0\}$, the
minimum over an empty set being $+\infty$; both are positive. Ties have probability zero.
On $\{z=k,h=k\}$ the residual is $r=\sigma_k|e|$.

\emph{Composition of $A_k$.} The mass of $A_k$ splits into the clean units of group $k$
with $h=k$, of mass $(1-\varepsilon)\pi_k(1-\eta_k)$, and the rest, of mass $\delta_kp_k$; the
rest consists of clean units of other groups, of mass at most
$(1-\varepsilon)\sum_j\pi_jP_0(h\ne j\mid z=j)=(1-\varepsilon)\eta$, and of contaminated
units. Since $\delta_k\le\nu_k<\tfrac12$, $p_k>0$ and $\eta_k<1$. The conditional distribution
function of $r$ on $A_k$ is therefore
\[
F_k(t)=(1-\delta_k)F^{\mathrm{own}}_k(t)+\delta_kF^{\mathrm{oth}}_k(t),\qquad
F^{\mathrm{own}}_k(t)=\frac{P_0(\sigma_k|e|\le t,\ h=k\mid z=k)}{1-\eta_k},
\]
with $F^{\mathrm{oth}}_k$ some distribution function. Since $e$ is independent of $(x,z)$,
$P_0(\sigma_k|e|\le t\mid z=k)=G_k(t)$, and $P_0(\sigma_k|e|\le t,\,h=k\mid z=k)$ lies between
$G_k(t)-\eta_k$ and $G_k(t)$. Hence
\begin{equation}
(1-\delta_k)\frac{G_k(t)-\eta_k}{1-\eta_k}\;\le\;F_k(t)\;\le\;
(1-\delta_k)\frac{G_k(t)}{1-\eta_k}+\delta_k .
\label{eq:adp-sandwich}
\end{equation}

\emph{(i) Uniqueness.} By the geometry above, $F^{\mathrm{own}}_k$ has density
\[
f^{\mathrm{own}}_k(t)=\frac{\varphi(t/\sigma_k)}{\sigma_k(1-\eta_k)}
\,\E_0\bigl[\one\{t<U(x)\}+\one\{t<L(x)\}\bigm|z=k\bigr],
\]
which is positive exactly when $P_0(\max\{L(x),U(x)\}>t\mid z=k)>0$; this property is
monotone in $t$, so $f^{\mathrm{own}}_k>0$ on $[0,T_k)$ and $f^{\mathrm{own}}_k=0$ on
$[T_k,\infty)$ for some $T_k\in(0,\infty]$. If the median of $F_k$ were not unique, $F_k$
would equal $\tfrac12$ on an interval of positive length. On $[0,T_k)$ the function $F_k$
is strictly increasing, so the interval lies in $[T_k,\infty)$, where
$F^{\mathrm{own}}_k=1$ and $F_k\ge1-\delta_k>\tfrac12$, a contradiction. The median is
positive because $F^{\mathrm{own}}_k(0)=0$ gives $F_k(0)\le\delta_k<\tfrac12$.

\emph{(i) Bounds.} Let $m_k$ be the median. If $F_k(t)\ge\tfrac12$ then, by the upper
bound in \eqref{eq:adp-sandwich}, $G_k(t)\ge u^{-}_k:=(\tfrac12-\delta_k)(1-\eta_k)/(1-\delta_k)$,
so $m_k\ge G_k^{-1}(u^{-}_k)$. If $G_k(t)\ge u^{+}_k:=\eta_k+(1-\eta_k)/(2(1-\delta_k))$ then,
by the lower bound, $F_k(t)\ge\tfrac12$, so $m_k\le G_k^{-1}(u^{+}_k)$; here $0<u^{-}_k$ and
$u^{+}_k<1$ because $\delta_k<\tfrac12$ and $\eta_k<1$. Now
$(\tfrac12-\delta_k)/(1-\delta_k)\ge\tfrac12-\delta_k$ and
$1/(2(1-\delta_k))\le\tfrac12+\delta_k$ for $\delta_k\le\tfrac12$, whence
$u^{-}_k\ge(\tfrac12-\delta_k)(1-\eta_k)\ge\tfrac12-\nu_k>0$ and
$u^{+}_k\le\eta_k+(1-\eta_k)(\tfrac12+\delta_k)\le\tfrac12+\nu_k<1$. Since
$G_k^{-1}(u)=\sigma_k\Phi^{-1}((1+u)/2)$ is increasing and $s_k=1.4826\,m_k$, the bounds
follow.

\emph{(ii).} Split $\kappa(\theta^{0})=P_0(r>cs_h)$ according to whether $h=z$. On
$\{h=z=k\}$, $r=\sigma_k|e|$ and $s_h=s_k$, so
$P_0(h=z,\ r>cs_h)=\sum_k\pi_kP_0(\sigma_k|e|>cs_k,\ h=k\mid z=k)$, which lies between
$\sum_k\pi_k\{G'_k-\eta_k\}$ and $\sum_k\pi_kG'_k$ with
$G'_k=\Pr(\sigma_k|e|>cs_k)=2\Phi(-cs_k/\sigma_k)$. The event $\{h\ne z\}$ contributes
between $0$ and $\eta$. Inserting the bounds of (i) on $s_k/\sigma_k$ into $G'_k$ gives
(ii); the expression for $\alpha(\theta^{0})$ is Theorem~\ref{thm:stepfixed}(ii), whose
Assumption~\ref{ass:reg} is verified in (iii) below, or directly from
\eqref{eq:adp-cleanonly}.

\emph{(iii) Nonsingularity.} By \eqref{eq:adp-cleanonly},
$M_k=(1-\varepsilon)\E_0[xx^{\top}\one_{W_k}]$, and $W_k\supseteq\{z=k,\,h=k,\,\sigma_k|e|\le cs_k\}$.
In the positive semidefinite order,
\begin{multline*}
\E_0\bigl[xx^{\top}\one\{z=k,\,h=k,\,\sigma_k|e|\le cs_k\}\bigr]\\
=\E_0\bigl[xx^{\top}\one\{z=k,\,\sigma_k|e|\le cs_k\}\bigr]
-\E_0\bigl[xx^{\top}\one\{z=k,\,h\ne k,\,\sigma_k|e|\le cs_k\}\bigr],
\end{multline*}
where the first term equals $\pi_kG_k(cs_k)\E_0[xx^{\top}\mid z=k]$ by the independence of
$e$ and $(x,z)$, and the second is at most $M_x^{2}\pi_k\eta_kI$ because
$xx^{\top}\preceq\|x\|^{2}I$. With $G_k(cs_k)=2\Phi(cs_k/\sigma_k)-1\ge2\Phi(cq(\tfrac12-\nu_k))-1$
by (i), $\lambda_{\min}(M_k)\ge(1-\varepsilon)\pi_k\lambda_k>0$. The other two parts of
Assumption~\ref{ass:reg} were shown in (i).

\emph{(iii) Bias.} $\theta^{\infty}_k-\theta^{0}_k=M_k^{-1}\E[x(y-x^{\top}\theta^{0}_k)\one_{W_k}]$,
and by \eqref{eq:adp-cleanonly} the expectation is $(1-\varepsilon)$ times its $P_0$-version.
Split the latter by $z$. On $\{z=k\}$, $y-x^{\top}\theta^{0}_k=\sigma_ke$ and
$W_k\cap\{z=k\}=\{z=k,\,h=k,\,|e|\le cs_k/\sigma_k\}$, so
\begin{multline*}
\E_0\bigl[xe\,\one\{z=k,\,h=k,\,|e|\le cs_k/\sigma_k\}\bigr]\\
=\E_0\bigl[x\one\{z=k\}\bigr]\,\E\bigl[e\one\{|e|\le cs_k/\sigma_k\}\bigr]
-\E_0\bigl[xe\,\one\{z=k,\,h\ne k,\,|e|\le cs_k/\sigma_k\}\bigr].
\end{multline*}
The first term vanishes by the symmetry of $e$; the second has norm at most
$(cs_k/\sigma_k)M_x\pi_k\eta_k$, so the contribution of $\{z=k\}$ is at most
$cs_kM_x\pi_k\eta_k$ in norm. On $\{z=j\}$, $j\ne k$, the factor $|y-x^{\top}\theta^{0}_k|=r\le cs_k$
on $W_k$, so the contribution is at most $cs_kM_xP_0(z=j,\,h=k)$ in norm, and summing over
$j\ne k$ gives at most $cs_kM_x\eta$. Altogether
$\|\E[x(y-x^{\top}\theta^{0}_k)\one_{W_k}]\|\le(1-\varepsilon)cs_kM_x(\pi_k\eta_k+\eta)\le2(1-\varepsilon)cs_kM_x\eta$,
and dividing by $\lambda_{\min}(M_k)\ge(1-\varepsilon)\pi_k\lambda_k$ gives the bound.

\emph{The special case.} If $\varepsilon=0$ and $\eta=0$ then $\delta_k=\nu_k=0$, so
$s_k=q(\tfrac12)\sigma_k$ by (i), $\kappa=2\Phi(-cq(\tfrac12))$ by (ii), and the bias bound
is zero.
\end{proof}

Three comments. First, the constant in the limit of the flagged fraction is
$2\Phi(-c\,s_k/\sigma_k)$ and not $2\Phi(-c)$. On the units assigned to surface $k$ the
minimum-over-surfaces residual is the residual from surface $k$ itself, so the per-group
median is a median of own-surface residuals; it is pulled down because the clean units of
group $k$ with the largest errors are nearer to another surface and are assigned there, and
pulled up by the units of other groups and by the contamination assigned to surface $k$.
Theorem~\ref{thm:steptrue}(i) bounds the net effect by $\nu_k=\delta_k+\eta_k/2$, which is
of the order of $\eta+\varepsilon$. The first effect is intrinsic to any scale computed
within an assigned group and vanishes only when the groups do not overlap. Second, the refit
is exactly unbiased at $\theta^{0}$ when the groups do not overlap and moves by $O(\eta)$
when they do; the source is the asymmetric truncation of the error on $\{h=k\}$, not the
contamination, which the step removes entirely under Assumption~\ref{ass:far}. Third,
Algorithm~\ref{alg:adaptive} reports the flagged fraction recomputed at the refitted
parameter, that is $\hat\alpha_n(\hat\theta)$ rather than $\hat\alpha_n(\theta)$;
Theorem~\ref{thm:stepfixed} concerns the flag set the refit uses. The limit of the reported
fraction is $\alpha(\theta^{\infty}(\theta))$ whenever $\alpha$ is continuous at
$\theta^{\infty}(\theta)$ and the convergence of $\hat\alpha_n$ is uniform near it, which we
do not prove.
